\chapter{Reactor Physics Fundamentals}
\label{ch:physics}

This chapter assembles the minimum reactor physics needed for a dynamic theory:
fission and the neutron balance, the multiplication factor and reactivity, the
six-factor formula, and the meaning of the neutron lifetime. Readers with a
background in nuclear engineering may skim it; the notation introduced here is
used everywhere later.

\section{Fission and the neutron economy}

When a $^{235}$U nucleus absorbs a thermal neutron it may fission, splitting into
two fragments and releasing on average $\nu\approx 2.43$ prompt neutrons together
with about $200\unit{MeV}$ of energy. The energy is distributed roughly as
$168\unit{MeV}$ in the kinetic energy of the fragments (deposited essentially at
the point of fission, in the fuel), with the remainder in prompt neutrons,
prompt gamma rays, and---importantly for later---the radioactive decay of the
fission products. Of the recoverable energy, about $93\%$ appears promptly and
about $7\%$ is released later as decay heat (Chapter~\ref{ch:decay}).

A self-sustaining chain reaction requires that, on average, exactly one of the
$\nu$ neutrons from each fission goes on to cause another fission. The other
neutrons are lost to non-fission capture (in fuel, moderator, coolant,
structure, and control materials) or to leakage out of the core. The bookkeeping
of these gains and losses is the neutron economy.

\section{The multiplication factor and reactivity}

Define the effective multiplication factor
\begin{equation}
\keff = \frac{\text{neutrons produced in one generation}}
             {\text{neutrons lost (absorbed + leaked) in the preceding generation}}.
\end{equation}
The three regimes are
\begin{equation}
\keff < 1 \ (\text{subcritical}),\qquad
\keff = 1 \ (\text{critical}),\qquad
\keff > 1 \ (\text{supercritical}).
\end{equation}
A critical reactor sustains a constant population; a supercritical one grows.

Because the interesting dynamics involve small departures from criticality, it is
convenient to measure the distance from $\keff=1$ directly. The
\textbf{reactivity} is
\begin{equation}
\boxed{\ \rho \equiv \frac{\keff-1}{\keff}\ }
\label{eq:reactivity}
\end{equation}
which is dimensionless, zero at criticality, positive above it, and negative
below. Reactivity is small in practice: a useful operating insertion might be a
few hundred pcm (recall $1\unit{pcm}=10^{-5}$), and one ``dollar'' equals
$\beta\approx 650\unit{pcm}$ for $^{235}$U. We will almost always work to first
order in $\rho$, where $\rho\approx \keff-1$.

\begin{definition}[Reactivity in dollars]
The reactivity expressed in units of the delayed-neutron fraction,
$\$\equiv \rho/\beta$, is the natural measure of transient severity. At
$\rho=\beta$ ($\$=1$) the reactor is prompt-critical: the prompt neutrons alone
sustain the chain, and delayed neutrons are no longer needed to maintain
growth.
\end{definition}

\section{The six-factor formula}

The effective multiplication factor for a finite thermal reactor factors into six
physically distinct processes:
\begin{equation}
\keff = \eta\, f\, p\, \varepsilon\, P_{\!f}\, P_{\!t},
\label{eq:sixfactor}
\end{equation}
where
\begin{itemize}[nosep]
\item $\eta$ is the reproduction factor (neutrons produced per thermal neutron
absorbed in fuel);
\item $f$ is the thermal utilisation (fraction of thermal absorptions occurring
in fuel rather than in moderator, coolant, or structure);
\item $p$ is the resonance escape probability (fraction of neutrons that slow
past the $^{238}$U capture resonances without being absorbed);
\item $\varepsilon$ is the fast-fission factor (a small boost from fissions
caused by fast neutrons);
\item $P_{\!f}$ and $P_{\!t}$ are the fast and thermal non-leakage
probabilities.
\end{itemize}
The product $\eta f p \varepsilon$ is the infinite-medium multiplication factor
$k_\infty$, and $\keff = k_\infty P_{\!f} P_{\!t}$.

The value of the six-factor formula for our purposes is that it localises
feedback. Temperature and voids do not change $\keff$ abstractly; they change
\emph{specific factors}. The Doppler effect (Chapter~\ref{ch:feedback}) acts on
the resonance escape probability $p$ through the broadening of the $^{238}$U
resonances. Moderator heating and voiding act on $f$ and $p$ through changes in
moderator density and the moderation ratio. Expansion changes leakage,
$P_{\!f}P_{\!t}$. Tracking \emph{which} factor a given feedback alters, and with
\emph{which sign}, is the discipline that distinguishes a stable lattice from an
unstable one.

\keypoint{Whether a reactor is under- or over-moderated is, in six-factor terms,
a competition: voiding the coolant tends to raise $p$ (less parasitic capture
and resonance absorption) but lower $f$ and worsen moderation. In an
under-moderated lattice the harmful effects win and reactivity falls; in an
over-moderated lattice (like the RBMK with its graphite) the helpful effect on
$p$ and the loss of the water absorber win, and reactivity rises.}

\section{The neutron lifetime and generation time}

Two closely related timescales govern how fast the population responds.

The \textbf{prompt neutron lifetime} $\ell$ is the mean time from a neutron's
birth to its absorption or leakage. In a thermal reactor it is dominated by the
slowing-down time plus the thermal diffusion time and is of order
$10^{-4}\unit{s}$ in a light-water reactor, somewhat longer in a graphite system.

The \textbf{prompt neutron generation time} is
\begin{equation}
\Lambda = \frac{\ell}{\keff}\approx \ell \quad(\keff\approx 1),
\end{equation}
the mean time between successive neutron generations. We use $\Lambda$ in the
kinetics equations; for $^{235}$U thermal systems a representative value is
$\Lambda\sim 10^{-5}\!-\!10^{-4}\unit{s}$. Its smallness is precisely why prompt
supercriticality is so dangerous: with $\Lambda=10^{-5}\unit{s}$ and a prompt
reactivity of just $0.2\%$, the prompt growth rate is
$\rho_{\text{prompt}}/\Lambda\sim 200\unit{s^{-1}}$, an e-folding every
$5\unit{ms}$.

\section{Neutron flux, cross-sections, and reaction rates}

The neutron flux $\phi(\mathbf r,E,t)$ is the product of neutron density and
speed; it measures the rate at which neutrons traverse unit area. Reaction rates
are products of flux and macroscopic cross-section,
\begin{equation}
R_x = \Sigma_x \phi, \qquad \Sigma_x = N\sigma_x,
\end{equation}
where $N$ is the number density of nuclei and $\sigma_x$ the microscopic
cross-section for reaction $x$ (fission, capture, scattering). Cross-sections
depend strongly on neutron energy; the resonance structure of $^{238}$U capture
near a few eV, and the $1/v$ behaviour of many thermal absorbers, are central to
feedback and to xenon poisoning respectively.

The fission power density is
\begin{equation}
q'''(\mathbf r,t) = E_f\, \Sigma_f\, \phi(\mathbf r,t),
\end{equation}
with $E_f\approx 200\unit{MeV}$ the recoverable energy per fission. Integrated
over the core, this is the thermal power that the kinetics equations track and
that the thermal-hydraulic system must remove.

\section{From the transport equation to point kinetics}

The full description of the neutron field is the time-dependent neutron transport
(or, in diffusion approximation, the neutron diffusion) equation, a
seven-dimensional integro-differential equation in space, energy, angle, and
time. For stability analysis we usually do not need its full spatial detail.
If the spatial shape of the flux changes slowly compared with its amplitude---the
\emph{point} approximation---we can separate variables,
\begin{equation}
\phi(\mathbf r,E,t)\approx \psi(\mathbf r,E)\, n(t),
\end{equation}
and reduce the dynamics to ordinary differential equations for the amplitude
$n(t)$ and the delayed-neutron precursors. This reduction yields the point
reactor kinetics equations of Chapter~\ref{ch:pk}. The approximation is excellent
for compact cores responding coherently, and it is the workhorse of reactor
dynamics. Its principal failure---when different regions of a large core behave
differently---gives rise to the spatial instabilities of
Chapter~\ref{ch:xenon}, and was a real complicating factor in the large RBMK core
at Chernobyl, where the point picture understates the danger of an asymmetric,
axially-peaked flux.


\section*{Exercises}
\addcontentsline{toc}{section}{Exercises}
\begin{enumerate}[leftmargin=*,itemsep=2pt]
\item Starting from $\keff=1.002$, compute the reactivity $\rho$ in pcm and in dollars (take $\beta=0.0065$).
\item Using the six-factor formula, identify which factor the Doppler effect modifies and with what sign.
\item Estimate the prompt growth rate $\rho_{\text{prompt}}/\Lambda$ for $\rho_{\text{prompt}}=0.003$ and $\Lambda=10^{-5}\,$s, and the corresponding e-folding time.
\item Explain physically why voiding an over-moderated lattice raises the resonance escape probability $p$.
\item Why does the point-kinetics approximation fail for large cores, and what phenomenon (Chapter 10) results?
\end{enumerate}
